\documentclass[12pt,a4paper]{article}
% ── Font encoding (needed for accented characters: ą, ń, ż, etc.) ────────────
\usepackage[T1]{fontenc}

% ── Mathematics ───────────────────────────────────────────────────────────────
\usepackage{amsmath,amssymb,amsthm}

% ── Figures ───────────────────────────────────────────────────────────────────
\usepackage{graphicx}
\usepackage{float}
\usepackage[caption=false,font=footnotesize]{subfig}

% ── Tables ────────────────────────────────────────────────────────────────────
\usepackage{array}

% ── TikZ ──────────────────────────────────────────────────────────────────────
\usepackage{tikz}
\usetikzlibrary{shapes.geometric,arrows.meta,positioning,calc}

% ── Float control ─────────────────────────────────────────────────────────────
\usepackage{placeins}
\usepackage{stfloats}   % allows table*/figure* at bottom of two-column pages

% Relax float placement so pages fill before floats are deferred.
% Defaults (0.7/0.3/0.2/0.5) are too strict for a float-heavy two-column paper
% and cause large blocks of white space.
\renewcommand{\topfraction}{0.9}         % max page fraction for top floats
\renewcommand{\bottomfraction}{0.9}      % max page fraction for bottom floats
\renewcommand{\textfraction}{0.07}       % min page fraction that must be text
\renewcommand{\floatpagefraction}{0.75}  % min fill before a float-only page
\setcounter{topnumber}{3}
\setcounter{bottomnumber}{2}
\setcounter{totalnumber}{5}

% ── Pseudocode ────────────────────────────────────────────────────────────────
\usepackage{algorithm}
\usepackage{algpseudocode}
\algrenewcommand\algorithmicrequire{\textbf{Input:}}
\algrenewcommand\algorithmicensure{\textbf{Output:}}
\algrenewcommand\alglinenumber[1]{{\footnotesize#1}}

% ── Custom notation ───────────────────────────────────────────────────────────
\newcommand{\alphah}{\hat{\alpha}}          % estimated alpha
\newcommand{\Ph}{\hat{P}}                   % empirical PDF
\newcommand{\rhoh}{\hat{\rho}}              % estimated autocorrelation
\newcommand{\VRh}{\widehat{VR}}             % estimated variance ratio
\newcommand{\Etau}{\mathbb{E}}              % expectation

\addbibresource{references.bib}

\onehalfspacing

% ─── Title ────────────────────────────────────────────────────────────────────
\title{\textbf{Literature Review}\\[0.4em]
       \large Empirical Lévy Scaling in Short-Horizon Futures Returns:\\
       Estimation, Sensitivity, and Autocorrelation Bias}
\date{\today}

\begin{document}
\maketitle
\tableofcontents
\newpage

% ══════════════════════════════════════════════════════════════════════════════
\section{Introduction}
% ══════════════════════════════════════════════════════════════════════════════

This review surveys the literature most directly relevant to a study of the
probability density function (PDF) of futures price increments, with four
principal aims: (i)~estimating the Lévy stability exponent $\alpha$ from
peak-scaling regressions and PDF collapses; (ii)~testing the robustness of
$\hat\alpha$ to methodological choices including de-seasonalization and
estimator selection; (iii)~examining the asymmetry of the empirical increment
distribution; and (iv)~deriving and validating a bias correction for the
peak-scaling estimator under short-term mean reversion.

The literature is organized into six thematic areas: the foundational theory
of Lévy-stable distributions and their financial motivation
(Section~\ref{sec:stable}); the seminal empirical evidence for heavy-tailed
and Lévy-like scaling in financial markets (Section~\ref{sec:empirical});
the estimation of tail and scaling exponents, including finite-sample cautions
(Section~\ref{sec:estimation}); the variance-ratio literature on short-term
dependence (Section~\ref{sec:vr}); intraday volatility seasonality and its
removal (Section~\ref{sec:seasonality}); and structure-function and
multifractal extensions of simple Lévy scaling (Section~\ref{sec:structure}).
A brief concluding section positions the present study relative to each thread
(Section~\ref{sec:position}).

Figure~\ref{fig:pipeline} gives a schematic overview of the analytical
pipeline, showing how the six sections of the empirical analysis relate to
one another.

Throughout this review, ``Lévy stable'' refers to the class of stable
distributions with characteristic exponent $\alpha \in (0,2]$, $\alpha = 2$
being the Gaussian special case. We follow the sign convention of
\textcite{SamorodTaqqu1994}, under which the characteristic function of a
symmetric stable variable $X \sim S_\alpha(\gamma,0,0)$ is
$\mathbb{E}[e^{iqX}] = \exp(-\gamma|q|^\alpha)$.


% ══════════════════════════════════════════════════════════════════════════════
\section{Lévy-Stable Distributions: Theory and Financial Motivation}
\label{sec:stable}
% ══════════════════════════════════════════════════════════════════════════════

\subsection{The stable class and the generalized central limit theorem}

The Gaussian distribution occupies a privileged role in probability theory as
the unique non-degenerate limit law of normalized sums of independent,
identically distributed (i.i.d.) random variables with finite variance,
a result known as the classical central limit theorem (CLT). The natural generalization, due to
Lévy and Khintchine, establishes that the class of \textit{stable
distributions} is the complete set of non-trivial distributional limits of
normalized partial sums of i.i.d.\ sequences
\parencite{SamorodTaqqu1994,GnedenkoKolmogorov1954}. A distribution
$F_\alpha$ is stable if and only if, for all $n \geq 1$, the sum of $n$
independent copies of $X \sim F_\alpha$ can be written as $n^{1/\alpha} X +
d_n$ for some centering constant $d_n$. This self-similarity under
convolution is the defining algebraic property of the stable family.

For $\alpha < 2$, stable distributions possess algebraically heavy tails:
$P(|X| > u) \sim C_\alpha u^{-\alpha}$ as $u \to \infty$, so that all moments
of order $\geq \alpha$ are infinite. Comprehensive treatments of the
mathematical theory are provided by \textcite{SamorodTaqqu1994} and
\textcite{Zolotarev1986}; the monograph of \textcite{Uchaikin2003} covers
applications in physics and finance.

\subsection{Mandelbrot's stable hypothesis for asset prices}

\textcite{Mandelbrot1963} was the first to propose stable distributions as a
model for asset-price changes, studying daily cotton prices and finding that
the observed heavy tails and lack of a finite-variance diffusion limit were
inconsistent with the Gaussian random-walk model of \textcite{Bachelier1900}.
Mandelbrot argued that the empirical stability exponent $\alpha$ was
approximately $1.7$, well inside the heavy-tailed regime, and that this was
consistent with the self-similar aggregation property: the distribution of
weekly price changes is a rescaled version of the distribution of daily
changes. This long-range self-similarity is the central empirical prediction
of the symmetric Lévy model tested in our study.

The hypothesis was further developed by \textcite{Fama1965}, who examined
daily returns on Dow Jones components and found that stable distributions with
$\alpha < 2$ provided a better fit than the Gaussian. These early findings were
subsequently contested on the grounds that the estimated $\alpha$ often exceeds
2 when computed from the Hill or regression estimators on finite samples
\parencite{Weron2001}, a point we return to in Section~\ref{sec:estimation}.

\subsection{Self-similarity and the scaling form}

The self-similar structure of a symmetric $\alpha$-stable law implies that the
density $P_{\alpha,\gamma}(x,\tau)$ of the $\tau$-step sum of i.i.d.\ stable
increments satisfies
\begin{equation}
  P_{\alpha,\gamma}(x,\tau)
  = (\gamma\tau)^{-1/\alpha}
    F_\alpha\!\left(\frac{x}{(\gamma\tau)^{1/\alpha}}\right),
  \label{eq:scaling}
\end{equation}
where $F_\alpha$ is the standardized stable density and $\gamma$ is the scale
parameter \parencite{SamorodTaqqu1994}. Setting $x = 0$ and using
$F_\alpha(0) = \Gamma(1/\alpha)/(\pi\alpha)$ yields the \textit{peak-scaling
formula}
\begin{equation}
  P(0,\tau) = \frac{\Gamma(1/\alpha)}{\pi\alpha}\,(\gamma\tau)^{-1/\alpha},
  \label{eq:peak}
\end{equation}
which gives a direct route to estimating $\alpha$ from the log-log regression
of the empirical PDF at the origin against the lag $\tau$.

\subsection{The truncated Lévy flight}

The idealized $\alpha$-stable model ($\alpha < 2$) has infinite variance, which
creates a tension with the finite empirical moments of any finite price dataset.
\textcite{MantegnaStanley1994} introduced the \textit{truncated Lévy flight}
(TLF) as a resolution: the distribution follows the stable power law over
intermediate scales but is exponentially cut off in the far tail, producing a
finite variance while preserving Lévy-like body scaling over the observable
range. The TLF framework has become standard in empirical Lévy finance
\parencite{BouchaudPotters2003,Cont2001,Mantegna1991}, because it allows
finite-variance dependence tools (variance ratios, autocorrelations) to be
applied to the same data for which the Lévy body is claimed, a hybrid
strategy that underpins the bias-correction analysis of this paper.

A brief controversy arose when \textcite{CopsonSmith1995} and
\textcite{Mantegna1995reply} debated whether the S\&P 500 data in
\textcite{MantegnaStanley1995} were truly consistent with a stable law;
the consensus view \parencite[Ch.~6]{BouchaudPotters2003} is that the TLF
provides an adequate phenomenological description for the lag ranges relevant to
intraday and short-horizon trading.


% ══════════════════════════════════════════════════════════════════════════════
\section{Empirical Evidence for Heavy-Tailed Scaling in Financial Markets}
\label{sec:empirical}
% ══════════════════════════════════════════════════════════════════════════════

\subsection{Mantegna and Stanley: intraday S\&P 500 scaling}

The first systematic empirical study of Lévy-type scaling in intraday data is
\textcite{MantegnaStanley1994,MantegnaStanley1995}. Using high-frequency
S\&P 500 index data, they showed that the central body of the return
distribution follows a symmetric Lévy law with $\alpha \approx 1.4$ for time
scales spanning three orders of magnitude (1 minute to 1000 minutes), while the
far tails fall off faster than any power law, consistent with the TLF.
\textcite{MantegnaStanley1995} demonstrated that the PDF collapses onto a
universal curve under the rescaling of Equation~\eqref{eq:scaling}, providing
the first clean empirical illustration of the procedure central to this paper.

\subsection{Cross-market universality: the inverse cubic and the 1.7 regime}

\textcite{Gopikrishnan1998} and \textcite{Gopikrishnan1999} extended the
analysis to individual US stocks using 40 million intraday transactions. They
found a power-law tail $P(|r| > u) \sim u^{-\zeta}$ with $\zeta \approx 3$
(the ``inverse cubic law''), beyond the Lévy-stable regime ($\zeta < 2$ would
be required for stability). \textcite{Weron2001} showed, however, that a
true Lévy-stable distribution with $\alpha \approx 1.7$-$1.8$ would produce
apparent Hill-estimator values of $\zeta \approx 3$ on samples of realistic
size, owing to finite-sample bias. This reconciliation between the Lévy and
inverse-cubic pictures motivates the multi-estimator approach in our paper and
the conservative interpretation of $\hat\alpha$ as a body-scaling parameter
rather than a tail-index claim.

\textcite{Plerou1999} showed that the scaling exponent is consistent across
a large cross-section of US equities, supporting the universality picture.
\textcite{Gabaix2003} provided a theoretical foundation for the inverse-cubic
law based on the trading of large institutions.

More recently, \textcite{Kadioglu2021} applied return-distribution analysis to
a broad cross-section of equity indices and showed that the tail behavior of
intraday returns continues to be well-characterized by power-law decay, with
evidence for the Lévy-stable regime surviving in the central body even on
data from 2010-2019. \textcite{Masoliver2021} provided an updated analysis of
financial return distributions over the period 2017-2020, confirming that the
inverse-cubic power law remains an appropriate global reference for tails at
intraday horizons, even as the distribution of returns has become heavier-tailed
in the post-crisis period. \textcite{Jiang2019} found similar scaling behavior
in Chinese equity markets.

\subsection{Futures markets specifically}

While most of the econophysics literature focuses on equity index or FX
returns, several studies have examined Lévy-type scaling in futures.
\textcite{Bouchaud2002} documented heavy-tailed and approximately stable
distributions for stock-index futures, bond futures, and commodity futures,
finding exponents consistent with $\alpha \in (1.5, 2)$ across markets.
\textcite{Cont2001} surveyed stylized facts across equity, FX, and futures
instruments, noting that heavy tails and volatility clustering are robust
across all asset classes. The recent survey \textcite{ContMiami2023}
verified that Cont's 2001 stylized facts, including heavy tails and the
absence of significant linear autocorrelation in returns, continue to hold
for intraday equity data from 2018-2019.

\subsection{Stability of the scaling picture under market evolution}

A natural concern is whether the Lévy-scaling picture documented in the 1990s
and early 2000s remains valid as market microstructure has evolved, driven by
algorithmic and high-frequency trading. \textcite{Masoliver2021} addressed this
directly, finding that the central-body scaling is persistent, although the
cross-over to the cut-off tail occurs at smaller multiples of the standard
deviation than in earlier data. \textcite{ContMiami2023} confirmed the
stability of heavy-tail and autocorrelation stylized facts in 2018-2019 Dow
Jones constituent data. These findings justify applying the Lévy peak-scaling
framework to contemporary futures data.


% ══════════════════════════════════════════════════════════════════════════════
\section{Estimation of the Scaling and Tail Exponent}
\label{sec:estimation}
% ══════════════════════════════════════════════════════════════════════════════

\subsection{Log-log regression (peak-scaling) estimator}

The peak-scaling estimator, which regresses $\log\widehat{P}(0,\tau)$ on
$\log\tau$ and infers $\hat\alpha = -1/\hat\beta$, was introduced by
\textcite{MantegnaStanley1994}. It has the advantage of being entirely
nonparametric with respect to the stable density shape, depending only on the
$\tau^{-1/\alpha}$ scaling of the central value. However, because it uses only
the information at $x = 0$, it discards information in the body and tails,
making it less efficient than maximum-likelihood or characteristic-function
estimators.

\textcite{Weron2001} documented that log-log regression systematically
overestimates $\alpha$ for stable distributions with $\alpha$ near 2 on finite
samples, because the asymptotic tail has not yet emerged. This paper takes this
caution seriously by (i)~reporting $R^2$ as a diagnostic for the linearity
assumption, (ii)~comparing three independent estimators, and (iii)~providing
the bias-correction formula for mean reversion derived in Extension~3.

\subsection{Collapse-quality estimator}

The PDF collapse estimator, which chooses $\alpha$ to minimize the spread of the
rescaled empirical curves $\tau^{1/\alpha}\widehat{P}(\tau^{1/\alpha}\tilde{x},\tau)$
around a common function, was proposed as a data-driven alternative to
regression-based methods by \textcite{Drozdz2003}. A closely related
``curve collapse'' approach has been used extensively in condensed-matter and
statistical physics to identify critical exponents \parencite{Houdayer2004}.
In the financial context, the collapse minimizes a mean-squared log-deviation
score, which weights the central body and the tails more evenly than a
pure-origin regression. The method is described and applied by
\textcite{BouchaudPotters2003} (Ch.~2) as a consistency check on the
peak-scaling estimate.

\subsection{Fractional-order structure functions}

Structure functions $S_q(\tau) = \mathbb{E}[|\Delta P_\tau|^q]$ were introduced
to finance by \textcite{Mandelbrot1997} and \textcite{MuzyBacry1993} in the
context of testing for multifractal scaling. Under exact Lévy scaling,
$S_q(\tau) \propto \tau^{q/\alpha}$ for $q < \alpha$; the implied $\alpha$
estimates from different values of $q$ should agree. Departures from linearity
in the scaling exponent $\zeta(q) = q/\alpha$ signal multifractal or
slowly-varying corrections. \textcite{BouchaudPotters2003} (Ch.~8) provides a
comprehensive treatment. For the fractional moment orders $q < 1$ used in this
paper, the estimator is particularly robust to outliers and converges reliably
on samples of moderate size \parencite{Uchaikin2003}.

\textcite{Barunik2010} provided Monte-Carlo evidence comparing Hill, Pickands,
and moment-based estimators, confirming that for stable-distributed data with
$\alpha \in [1.5,2)$, log-log moment-based estimators (including fractional
structure functions) have lower bias than the Hill estimator.

\subsection{Maximum-likelihood and characteristic-function estimators}

For comparison, we note that maximum-likelihood estimators (MLE) for stable
distributions are available \parencite{Nolan2001} and have been shown to be
asymptotically efficient \parencite{DuMoutierde1983}. Characteristic-function
(CF) based estimators, including those of \textcite{Feuerverger1981} and
\textcite{Press1972}, exploit the analytic form of the CF. Both families are
computationally more demanding than the log-log regression and require a fully
parametric specification of the stable family, losing the model-free character
of the peak-scaling approach. The present study intentionally uses
semi-parametric methods (tip scaling, collapse, and structure functions) to
reduce sensitivity to misspecification of the tail shape.

\subsection{The finite-variance debate and the Hill estimator}

A persistent question in the empirical literature is whether stock and futures
returns have $\alpha < 2$ (infinite variance) or $\alpha > 2$ (finite variance,
power-law but non-stable tails). The Hill estimator consistently reports
$\hat\alpha \approx 3$ for large cross-sections of equities
\parencite{Gopikrishnan1999}, which would place the tail exponent outside the
Lévy-stable regime. \textcite{Weron2001} showed definitively that the Hill
estimator and log-log regression both overestimate $\alpha$ for stable
distributions near $\alpha = 2$: samples of $10^6$ or more are required before
the true tail emerges, while practical financial datasets rarely exceed $10^4$
to $10^5$ observations at intraday frequency.

The reconciliation proposed by \textcite{Weron2001} and adopted here is that
the body of the empirical distribution (where the bulk of the sample mass lies)
is consistent with a Lévy-stable law with $\alpha \approx 1.7$-$1.8$, while
the far tail appears steeper due to the truncation of real-world distributions.
The peak-scaling and collapse estimators are sensitive primarily to the central
body and are therefore appropriate tools for testing the TLF picture, even
though they would not reliably recover the true tail index of a pure stable law
on these sample sizes.


% ══════════════════════════════════════════════════════════════════════════════
\section{Short-Term Dependence and the Variance-Ratio Framework}
\label{sec:vr}
% ══════════════════════════════════════════════════════════════════════════════

\subsection{The random-walk hypothesis and the variance-ratio test}

The random-walk hypothesis (RWH) asserts that successive price increments are
independent, so the variance of the $\tau$-step change scales linearly with
$\tau$: $\mathrm{Var}(\Delta P_\tau) = \tau\,\mathrm{Var}(\delta_1)$.
\textcite{LoMacKinlay1988} introduced the variance-ratio test
\begin{equation}
  VR(\tau) = \frac{\mathrm{Var}(\Delta P_\tau)}{\tau\,\mathrm{Var}(\delta_1)},
  \label{eq:vr}
\end{equation}
which equals 1 under the RWH and deviates from 1 under mean reversion
($VR < 1$) or momentum ($VR > 1$). They derived asymptotically valid
$z$-statistics for both homoscedastic and heteroscedastic increments and
applied the test to weekly US stock returns, finding $VR > 1$ (positive serial
correlation) at short horizons. The comprehensive treatment of the test
appears in \textcite{CampbellLoMacKinlay1997}, Chapter~2.

The exact relationship between the variance ratio and the autocorrelation
function is
\begin{equation}
  VR(\tau)
  = 1 + 2\sum_{k=1}^{\tau-1}\!\left(1 - \frac{k}{\tau}\right)\rho(k),
  \label{eq:vr_rho}
\end{equation}
where $\rho(k) = \mathrm{Corr}(\delta_t, \delta_{t+k})$. This identity is the
basis for both of the variance-ratio estimators implemented in our code.

\subsection{Mean reversion at short horizons in equity and futures markets}

At the intraday and short-horizon frequencies relevant to this paper,
the dominant finding is \textit{negative} serial correlation (mean reversion)
in futures and equity returns, not the positive correlation found at weekly
horizons by \textcite{LoMacKinlay1988}. Several mechanisms have been proposed:

\textbf{Bid-ask bounce.} \textcite{Roll1984} showed that even in the absence
of true serial correlation, transactions alternating between bid and ask prices
generate spurious negative first-order autocorrelation. The magnitude is
$\hat\rho(1) \approx -s^2/4$ where $s$ is the effective bid-ask spread. The
bid-ask bounce contributes a mechanically-induced mean reversion at the 1-minute
frequency used in this paper.

\textbf{Order-book pressure and price reversion.} \textcite{Bouchaud2004}
and \textcite{Farmer2006} documented that large market orders move prices, but
this impact partially reverts as the order book replenishes. At 1-minute
resolution, this ``impact reversion'' contributes to the negative
autocorrelation of price changes.

\textbf{Empirical magnitude.} The most relevant evidence for futures markets
comes from \textcite{LoMacKinlay1990}, who found stronger mean reversion in
small-capitalization stocks. For index futures specifically,
\textcite{BouchaudPotters2003} (Ch.~6) report first-order autocorrelations
$\hat\rho(1) \approx -0.1$ to $-0.2$ at 1-minute frequency across major equity
index contracts, consistent with the values found in this paper.

\subsection{Implications for peak-scaling estimation}

The tension between the i.i.d.\ assumption of the Lévy scaling law and the
empirically documented mean reversion motivates the bias-correction analysis
of Extension~3 in this paper. Under weak dependence, the effective Lévy scale
parameter at lag $\tau$ is modified by the variance ratio:
$\gamma_{\mathrm{eff}}(\tau) = \gamma\tau \cdot VR(\tau)^{\alpha/2}$
(derived from the Fourier-inversion formula of Equation~\eqref{eq:peak} and
the second-cumulant correction of \textcite{BouchaudPotters2003}, Ch.~6).
For a mean-reverting process $VR(\tau) < 1$, so the effective scale is smaller
than the i.i.d.\ prediction, causing the raw peak-scaling regression to
overestimate $\alpha$ (bias toward the Gaussian boundary).

To our knowledge, this precise bias-correction formula for the peak-scaling
estimator under a truncated-Lévy body with mean-reversion has not been derived
and empirically validated in the prior literature. \textcite{BouchaudPotters2003}
discuss the qualitative direction of the effect, and
\textcite{SamorodTaqqu1994} (Section~7.2) derive the exact scale modification
for stable ARMA models, but the specific OLS-bias formula and its cross-market
validation are original contributions of the present study.

\subsection{Extensions and alternatives to the variance-ratio test}

\textcite{ChowDenning1993} extended the \textcite{LoMacKinlay1988} test to a
joint test over multiple lags using the studentized maximum modulus distribution.
\textcite{Wright2000} proposed rank and sign-based variance-ratio tests with
better finite-sample size. \textcite{Richardson1991} derived the exact
distribution of the variance-ratio statistic. These extensions are relevant for
inference on the variance ratio; in our context we use the variance ratio as an
input to a correction formula rather than as a test statistic, so we rely on
the point-estimator interpretation of \textcite{CampbellLoMacKinlay1997}.


% ══════════════════════════════════════════════════════════════════════════════
\section{Intraday Volatility Seasonality and De-Seasonalization}
\label{sec:seasonality}
% ══════════════════════════════════════════════════════════════════════════════

\subsection{The U-shaped intraday volatility pattern}

It is a well-documented stylized fact that intraday price volatility follows a
deterministic pattern across the trading session. For most equity and futures
markets, volatility is elevated near the open (driven by overnight information
arrival and the unwinding of market-open imbalances), depressed at midday, and
rises again near the close (driven by end-of-day portfolio rebalancing and
hedging). This ``U-shape'' was documented early by \textcite{WoodMcInishOrd1985}
for NYSE equities and has since been replicated across asset classes and
geographies.

\subsection{Andersen and Bollerslev's decomposition model}

The systematic econometric treatment of intraday periodicity is due to
\textcite{AndersenBollerslev1997,AndersenBollerslev1998}, who showed that the
dominant intraday volatility patterns in FX and equity markets can be factored
as a deterministic time-of-day component multiplied by a stochastic process:
$\sigma_t = \hat\sigma_{m(t)} \cdot \tilde\sigma_t$, where $m(t)$ is the
minute-of-session. Removing the deterministic component by normalizing
$\tilde\delta_t = \delta_t / \hat\sigma_{m(t)}$ substantially improves the fit
of GARCH-type models for the residual stochastic volatility.

\textcite{AndersenBollerslev1998} applied this decomposition to Deutsche
Mark-Dollar FX returns and showed that the periodicity-adjusted series has
substantially more stable autocorrelation properties. For equity index futures
specifically, \textcite{Ederington1993} documented U-shaped volatility patterns
in S\&P 500 futures and \textcite{Lee1994} showed that de-seasonalization is
important for accurate volatility forecasting.

\subsection{Effects on tail and scaling estimates}

The interaction between intraday seasonality and tail exponent estimation has
received less attention in the literature. If the intraday volatility profile
is not removed, the estimated $\tau$-step increment distribution mixes
low-volatility midday increments with high-volatility open and close increments.
At short lags (small $\tau$), this mixture does not break self-similarity
because the scaling formula is applied at fixed $\tau$; however, the effective
scale parameter $\gamma$ in Equation~\eqref{eq:scaling} is biased by the time
of day at which increments are drawn, potentially distorting the peak-scaling
regression across lags.

\textcite{BouchaudPotters2003} (Ch.~6) discuss the de-seasonalization step as
standard practice before applying scaling analyses to intraday data. The
question of whether de-seasonalization materially shifts $\hat\alpha$
or improves collapse quality is examined across all eight markets in our
study. This constitutes, to our knowledge, the most systematic multi-market
test of the effect of de-seasonalization on Lévy peak-scaling estimates in the
futures context.

\subsection{Recent evidence on intraday patterns}

\textcite{Kadioglu2023} applied recurrence quantification analysis to intraday
data and confirmed the persistence of U-shaped volatility patterns in modern
equity markets, noting their stability across market regimes. The volatility
forecasting literature continues to treat de-seasonalization as a prerequisite
for high-frequency modeling \parencite{ParkEomKim2022,LiangWan2023}.


% ══════════════════════════════════════════════════════════════════════════════
\section{Structure Functions, Multifractality, and Scaling Robustness}
\label{sec:structure}
% ══════════════════════════════════════════════════════════════════════════════

\subsection{Structure functions in turbulence and finance}

Fractional-order structure functions $S_q(\tau) = \mathbb{E}[|\Delta P_\tau|^q]$
originated in Kolmogorov's theory of turbulence as a tool for characterizing
the intermittency of velocity increments \parencite{Kolmogorov1941}. The same
framework was imported into finance by \textcite{MuzyBacry1993,MuzyBacry1994}
and \textcite{Mandelbrot1997} to test for simple versus multifractal scaling.
Under a pure $\alpha$-stable law, the scaling exponent function $\zeta(q) =
q/\alpha$ is linear in $q$ (monofractal scaling); any systematic concavity of
$\hat\zeta(q)$ as a function of $q$ signals multifractality, i.e.\ the
presence of multiple scaling exponents.

\subsection{Evidence for and against multifractality in financial data}

\textcite{Calvet1997} and \textcite{Mandelbrot1997} provided early evidence
for multifractal scaling in FX data, and the multifractal random walk model of
\textcite{BacryMuzy2003} has become influential in the volatility-modeling
literature. However, the evidence is contested: \textcite{Cont2001} notes that
apparent multifractality can arise from volatility clustering (GARCH effects)
in finite samples, and \textcite{Eisler2004} showed that multifractal estimates
from individual stocks are sensitive to the choice of $q$ values and sample length.

More recently, \textcite{Saadaoui2024} applied multifractal detrended
fluctuation analysis to cryptocurrency returns and found structured multifractal
scaling, though with substantial sensitivity to the estimation window.
\textcite{Morales2012} provided a detailed comparison of scaling exponent
estimators for financial time series and found that the structure-function
approach is reliable for $q \in (0,1)$ but sensitive to outliers for $q > 1$.

The structure-function analysis in our paper uses $q \in \{0.25, 0.5, 0.75,
1.0, 1.25, 1.5\}$, spanning both the sub-linear and super-linear regimes, and
examines linearity of $\hat\zeta(q)$ as a consistency check on the Lévy
scaling picture rather than a test for multifractality per se.

\subsection{Asymmetry of return distributions}

The symmetric Lévy model assumes that the increment distribution is symmetric
around zero. A growing body of evidence suggests that this symmetry is
approximate but not exact in equity and futures markets. Negative skewness
(``leverage effect'') has been documented at daily and longer horizons
\parencite{Cont2001,Bouchaud2001}. At intraday horizons the direction of
asymmetry is less settled, but \textcite{Ait-Sahalia2002} found statistically
significant skewness in high-frequency equity returns, while
\textcite{BouchaudPotters2003} (Ch.~3) argue that the skewness at 1-minute
frequency is small enough that the symmetric Lévy approximation is adequate as
a first-order description.

The asymmetry analysis in Extension~2 of our paper addresses this question
directly: we test whether the positive and negative sides of the empirical PDF
are significantly different and whether symmetrization materially improves
collapse quality. This complements the theoretical asymmetric stable-distribution
literature \parencite{Zolotarev1986,Nolan2001}.

\subsection{Robustness analysis and sensitivity to methodological choices}

The question of how sensitive estimated scaling exponents are to modeling
and methodological choices has received growing attention. \textcite{Cont2001}
emphasized that different estimators can give qualitatively different answers
for the same dataset. \textcite{Weron2001} provided Monte Carlo evidence on the
bias and variance of common estimators. \textcite{Barunik2010} compared
tail-index estimators systematically and showed that Hill-based methods
are unreliable for $\alpha$ near 2.

Our robustness matrix (Extension~1) provides a systematic multi-estimator,
multi-preprocessing, multi-market sensitivity table in the spirit of
\textcite{Sala-i-Martin1997}'s sensitivity analysis methodology. To our
knowledge, no prior study has reported such a comprehensive sensitivity analysis
specifically for the Lévy peak-scaling estimate across a cross-section of
futures markets.


% ══════════════════════════════════════════════════════════════════════════════
\section{Positioning of the Present Study}
\label{sec:position}
% ══════════════════════════════════════════════════════════════════════════════

This paper contributes to each of the six literatures reviewed above.

\textbf{Relative to the stable-distribution literature.} We apply the
TLF framework of \textcite{MantegnaStanley1994,MantegnaStanley1995} to a
cross-section of eight futures markets, extending the single-instrument analyses
of the foundational econophysics literature to a broader empirical base.

\textbf{Relative to the empirical scaling literature.} Whereas most
studies focus on equity indices or individual stocks, we apply the Lévy
peak-scaling and PDF collapse approach directly to liquid futures markets
(equity index, fixed income, and commodity), testing the universality of
the scaling exponent $\alpha \approx 1.7$-$1.8$ reported in the literature.

\textbf{Relative to the estimation literature.} We compare three
independent estimators (tip scaling, collapse quality, fractional structure
functions) and construct a robustness matrix across all estimator-by-preprocessing
combinations. The comparison addresses the cautions of \textcite{Weron2001}
and \textcite{Barunik2010} about relying on a single estimator.

\textbf{Relative to the variance-ratio literature.} We use the
Campbell-Lo-MacKinlay variance ratio not as a test statistic but as a
correction factor in a peak-scaling bias formula, a novel application that
connects the dependence literature to the scaling literature. The resulting
corrected estimator $\hat\alpha_{\mathrm{corr}}$ is, to our knowledge, the
first closed-form bias correction for the Lévy peak-scaling estimator under
short-term mean reversion.

\textbf{Relative to the seasonality literature.} We provide the most
systematic multi-market test of the effect of de-seasonalization on Lévy
peak-scaling estimates, complementing the single-market studies of the
intraday volatility literature.

\textbf{Relative to the structure-function literature.} Our fractional
structure-function analysis serves as a cross-check on the Lévy
self-similarity assumption rather than a test for multifractality, consistent
with the interpretation in \textcite{BouchaudPotters2003}.


% ══════════════════════════════════════════════════════════════════════════════
\printbibliography[heading=bibintoc,title={References}]
% ══════════════════════════════════════════════════════════════════════════════

\end{document}
